% Partie : théorie de la démonstration
% Chapitre : calcul des séquents
% Section : règles de mG3i

\documentclass[../../../include/open-logic-section]{subfiles}

\begin{document}

\begin{table}
\begin{tabular}{@{}c@{}c@{}}
\multicolumn{2}{@{}c@{}}{Axiomes}\\[2ex]
$!C, \Gamma \Sequent \Delta, !C$ & $\lfalse, \Gamma \Sequent \Delta$
\\[2ex]
\multicolumn{2}{@{}c@{}}{Règles logiques}\\[2ex]
\Axiom$\lnot !A, \Gamma \fCenter \Delta, !A $
\RightLabel{\LeftR{\lnot}}
\UnaryInf$\lnot !A, \Gamma \fCenter \Delta$
\DisplayProof
&
\Axiom$!A, \Gamma \fCenter$
\RightLabel{\RightR{\lnot}}
\UnaryInf$ \Gamma \fCenter \Delta, \lnot !A $
\DisplayProof
\\[4ex]
\Axiom$ !A, !B, \Gamma \fCenter \Delta$
\RightLabel{\LeftR{\land}}
\UnaryInf$ !A \land !B, \Gamma \fCenter \Delta$
\DisplayProof
&
\Axiom$\Gamma \fCenter \Delta, !A$
\Axiom$ \Gamma \fCenter \Delta, !B$
\RightLabel{\RightR{\land}}
\BinaryInf$ \Gamma \fCenter \Delta, !A \land !B $
\DisplayProof
\\[4ex]\Axiom$!A, \Gamma \fCenter \Delta$
\Axiom$!B, \Gamma \fCenter \Delta$
\RightLabel{\LeftR{\lor}}
\BinaryInf$!A \lor !B, \Gamma \fCenter \Delta$
\DisplayProof
&
\Axiom$\Gamma \fCenter \Delta, !A, !B$
\RightLabel{\RightR{\lor}}
\UnaryInf$ \Gamma \fCenter \Delta, !A \lor !B$
\DisplayProof
\\[4ex]
\Axiom$!A \lif !B, \Gamma \fCenter \Delta, !A$
\Axiom$ !B, \Gamma \fCenter \Delta$
\RightLabel{\LeftR{\lif}}
\BinaryInf$!A \lif !B, \Gamma \fCenter \Delta$
\DisplayProof
&
\Axiom$ !A, \Gamma \fCenter !B$
\RightLabel{\RightR{\lif}}
\UnaryInf$ \Gamma \fCenter \Delta, !A \lif !B $
\DisplayProof
\\[4ex]
\Axiom$ !A(t), \lforall[x][!A(x)], \Gamma \fCenter \Delta$
\RightLabel{\LeftR{\lforall}}
\UnaryInf$ \lforall[x][!A(x)],\Gamma \fCenter \Delta$
\DisplayProof
&
\Axiom$ \Gamma \fCenter !A(a) $
\RightLabel{\RightR{\lforall}}
\UnaryInf$ \Gamma \fCenter \lforall[x][!A(x)]$
\DisplayProof
\\[4ex]
\Axiom$ !A(a), \Gamma \fCenter \Delta $
\RightLabel{\LeftR{\lexists}}
\UnaryInf$ \lexists[x][!A(x)], \Gamma \fCenter \Delta$
\DisplayProof
&
\Axiom$ \Gamma \fCenter \Delta, \lexists[x][!A(x)], !A(t) $
\RightLabel{\RightR{\lexists}}
\UnaryInf$ \Gamma \fCenter \Delta, \lexists[x][!A(x)]$
\DisplayProof
\end{tabular}
\caption{Règles de \Log{mG3i}, calcul intuitionniste à conséquent
multiple. Dans $\RightR{\forall}$ et $\LeftR{\exists}$, la
constante~$a$ doit être absente de la conclusion de l'inférence.
$\Gamma$ et $\Delta$ sont des multiensembles finis de formules.
Les termes d'instanciation sont clos.
Dans les axiomes d'identité, la formule~$!C$ doit être atomique.
Le texte original mentionne aussi \Log{mG1m}, obtenu à partir de
\Log{mG1i} en retirant les axiomes $\lfalse, \Gamma \Sequent \Delta$.}
%
%
\ollabel{tab:mG3i}
\end{table}

\end{document}
